\section*{Sammanfattning}\label{sec:sammanfattning}

\setlength{\emergencystretch}{2em}

Arbetsflöden för sakernas internet (IoT) kan koppla sensorhändelser till fysiska
eller simulerade åtgärder, men språkmodellsförslag är inte nödvändigtvis
kontraktsenliga eller godkända för utförande. Syftet var att utforma,
implementera och utvärdera Obids bidrag: ett avgränsat agentbaserat besluts- och
tillförlitlighetslager med en enda agent som utökar den ärvda
Yacoub-kompatibla n8n-gränsen mellan arbetsflöde och åtgärd och bevarar det
gemensamma åtgärdsgränssnittet. Lagret kombinerade två verktyg med enbart
läsåtkomst, begränsat minne, körtidsvalidering, deterministisk policy och en
exekverbar mekanism för människa i loopen (HITL). \texttt{CONFIG-OBID}
jämfördes med den ärvda baslinjen med en minimal agent,
\texttt{CONFIG-BASELINE}, under ett fryst protokoll med upprepade försök i ett
temperaturstyrt fläktscenario; alla planerade primära försök, inklusive fel och
avvikelser, behölls.

Totalt behölls 85 primära försök. \texttt{CONFIG-OBID} gav det förväntade
utfallet i 35/35 RQ1-försök. Körtidsvalideringen blockerade alla fem injicerade
ogiltiga åtgärder; i alla tio väntande HITL-försök hölls åtgärden kvar utan
frisläppning. Tilldelade godkännanden gav korrekt utfall i 5/5 försök och
tilldelade avslag i 4/5, eftersom en människa faktiskt godkände ett tilldelat
avslagsförsök; detta kvarstod som en protokollavvikelse. Ingen ogiltig eller
icke godkänd åtgärd passerade det gemensamma åtgärdsgränssnittet i de
observerade försöken. Jämfört med \texttt{CONFIG-BASELINE} skilde sig
tillförlitligheten vid hantering av felaktigt strukturerade indata och
undertryckning av duplicerade åtgärder; \texttt{CONFIG-OBID} uppvisade högre
medianlatens för automatiserad behandling i alla tre uppmätta fallfamiljer.
Resultaten visar att den avgränsade utökningen förenade förväntat
beslutsbeteende med exekverbara frisläppningskontroller i den utvärderade
miljön, men fastställer inte universell tillförlitlighet eller
produktionssäkerhet. Underlaget begränsas till en modell, ett temperaturstyrt
fläktscenario, fem repetitioner per cell och en simulerad fläkt vid
systemgränsen.

Nyckelord: n8n, agentbaserade arbetsflöden, sakernas internet,
körtidsvalidering, människa i loopen, arbetsflödesautomatisering

\setlength{\emergencystretch}{0pt}
